% !TEX program = xelatex
% LSPV-3 REPORT::MINFRONT1::2026-10-12T05:32:10Z
\documentclass[11pt]{article}
\usepackage[margin=1in]{geometry}
\usepackage{amsmath,amssymb,booktabs,longtable,array,enumitem,xcolor,hyperref}
\hypersetup{colorlinks=true,linkcolor=blue,urlcolor=blue}
\newcommand{\OPEN}{\textbf{OPEN}}
\newcommand{\proved}{\textbf{PROVED/EXACT}}
\newcommand{\known}{\textbf{KNOWN THEOREM}}
\newcommand{\comp}{\textbf{COMPUTATIONAL}}
\newcommand{\cand}{\textbf{CANDIDATE}}
\newcommand{\retired}{\textbf{FALSIFIED/RETIRED}}
\newcommand{\same}{\textsc{same\_realization}}

\title{LSPV-3 MINFRONT1: an exponential minimal-incompatibility frontier with a linear evaluator}
\author{UNIVERSAL\_CORE\_REDTEAM}
\date{12 October 2026}

\begin{document}
\maketitle

\begin{center}
\fbox{\parbox{0.94\linewidth}{
\textbf{Executive status.} P versus NP remains \OPEN.  This report proves an
exact structural counterexample: a proper duplicate-free 3-CNF family has
$2^{n-1}$ minimal incompatible source partial assignments, but its projected
compatibility predicate is evaluated by one linear scan with logarithmic
address space.  Therefore semantic frontier cardinality, even for an antichain
of minimal obstructions, does not by itself force materialization, extraction,
communication, substantial workspace, or superpolynomial time in an arbitrary
one-bit SAT computation.  The result is an easy-island falsification only; by
QGUARD1 it does not refute the surviving full-domain worst-case transfer.}}
\end{center}

\section{Target and exact quantifiers}

Let $X$ be the encoded 3-CNF domain and let $\mathcal D$ be deterministic
uniform exact SAT deciders.  The attacked claim was the unqualified promotion
schema
\begin{equation}
\tag{MIFO1}
\forall A\in\mathcal D\;\forall F\in\mathcal C:\quad
|\operatorname{MinBad}(H_F)|\ge 2^{\Omega(n(F))}
\Longrightarrow A(F)\text{ exposes or pays }2^{\Omega(n(F))}.
\end{equation}
Here ``pays'' was intended to cover time, writable memory, ordinary external
communication, explicit extraction, or reconstruction.  Without a separately
proved execution relation, this implication is false.

For a 3-CNF $F(X,Z)$ with declared source variables $X=(x_1,\ldots,x_n)$ and
auxiliaries $Z$, define the projected compatibility predicate
\[
H_F(\alpha)=1\quad\Longleftrightarrow\quad
\exists x\succeq\alpha\ \exists z\;F(x,z)=1,
\qquad \alpha\in\{0,1,\star\}^n.
\]
Order partial assignments by extension.  A minimal incompatible assignment is
an $\alpha$ with $H_F(\alpha)=0$ such that every strict weakening obtained by
replacing an assigned bit by $\star$ has compatibility one.

The model is an ordinary multi-tape deterministic machine receiving an indexed
3-CNF bit string of length $N$.  The source-variable declaration is part of the
family specification, not hidden advice.  Ordinary SAT output is one bit.  All
preprocessing, clause generation, clause checks, evaluator calls, extraction,
reconstruction and communication are separately metered in the audit.

\section{The parity-source family}

Fix $n\ge1$ and $b\in\{0,1\}$.  Let $P_{n,b}^{\rho}$ be the proper
duplicate-free 3-CNF obtained from a left-chain or balanced binary XOR circuit
for
\[
x_1\oplus\cdots\oplus x_n=b.
\]
Each gate $z=a\oplus c$ is encoded by the four clauses
\[
(a\vee c\vee\neg z),\quad(a\vee\neg c\vee z),\quad
(\neg a\vee c\vee z),\quad(\neg a\vee\neg c\vee\neg z).
\]
The final output literal $\ell$ is forced true, without repeated literals, by
fresh variables $p,q$ and
\[
(\ell\vee p\vee q),\ (\ell\vee p\vee\neg q),\
(\ell\vee\neg p\vee q),\ (\ell\vee\neg p\vee\neg q).
\]
For $n\ge2$ there are $4n$ clauses and $2n+1$ variables.  Indexed-literal
encoding has $N=\Theta(n\log n)$ bits.  Bags equal to the three variables of
each XOR gate, with the force bag attached to the output gate, give a width-two
primal tree decomposition for either layout.

\subsection{Exact theorem}

\paragraph{PARITY-PROJECTED-FRONTIER1 (\proved, pending independent review).}
For every $n\ge1$, $b\in\{0,1\}$ and either layout $\rho$,
\begin{align*}
H_{P_{n,b}^{\rho}}(\alpha)=1
&\Longleftrightarrow
\bigl(\exists i:\alpha_i=\star\bigr)
\ \vee\
\bigl(\alpha\in\{0,1\}^n\land\bigoplus_i\alpha_i=b\bigr),\\
|\operatorname{MinBad}(H_{P_{n,b}^{\rho}})|&=2^{n-1}.
\end{align*}
Every member of $\operatorname{MinBad}$ is a full wrong-parity assignment, and
these assignments form an antichain.

\paragraph{Proof.}
If a source coordinate is unassigned, choose all other unassigned coordinates
arbitrarily and choose the last one to achieve target parity $b$.  The XOR-gate
values are then uniquely extended and the force clauses are satisfiable.  If
the source is fully assigned, the gate clauses and final force are satisfiable
exactly when its parity is $b$.  Exactly half of the $2^n$ full assignments have
wrong parity.  Weakening any one of them introduces a star and therefore
restores compatibility.  Hence all and only those $2^{n-1}$ assignments are
minimal incompatible.  The displayed evaluator scans once for a star and,
only if none exists, accumulates parity.  It runs in $O(n)$ time and
$O(\log n)$ address space (one parity bit plus scan index). \hfill$\square$

\paragraph{MINIMAL-INCOMPATIBILITY-FRONTIER-OBLIGATION1 (\retired).}
The family has an exponentially large minimal-incompatibility antichain but a
compact direct evaluator.  Its SAT answer is always YES after canonical-family
recognition.  Thus no conclusion about materializing the frontier, executing
one operation per frontier element, communicating it, or using exponential
time/space follows from its cardinality.

\section{Arbitrary deciders, wrappers and the quantifier guard}

Under the counterfactual assumption $\mathrm{SAT}\in\mathrm P$, let $A$ be any
polynomial SAT decider.  Define the uniform recognize-first wrapper
\[
W_A(F)=\begin{cases}
1,&F\text{ is a canonical }P_{n,b}^{\rho},\\
A(F),&\text{otherwise.}
\end{cases}
\]
Canonical recognition and parity-family validation are polynomial in encoded
length.  Hence $W_A$ is a polynomial exact SAT decider that need not enumerate
or expose $\operatorname{MinBad}$ on these inputs.  COPY, DELAY, RECOMP,
clause-order, literal-order, redundant-clause, interpreter and DEAD wrappers
further show that named internal events are not representation-invariant.

This does \emph{not} negate
\[
\forall A\in\mathcal D\ \exists F\in X:\ R(A,F),
\]
because the counterexample is an easy island.  QGUARD1 proves that avoiding an
obligation on a proper island is compatible with a worst-case obligation on
the full language.  The surviving repair must therefore be full-domain and
existential-in-input, or provide a reduction mapping every hard source to the
charged family while preserving the claimed execution obligation.

\section{Reproducible audit}

Program:
\path{experiments/review/minimal_incompatibility_frontier_audit.py}.
Result:
\path{results/raw/review/2026-10-12_minimal_incompatibility_frontier.json}.
The invocation was
\begin{quote}\small
\texttt{python minimal\_incompatibility\_frontier\_audit.py
--max-exhaustive-n 8 --seed 20261012 --output <result>}.\normalsize
\end{quote}
It completed in $2.4813$ seconds.  Exhaustive checks covered $n=1,\ldots,8$,
both parity targets, both chain layouts, all $3^n$ source partial assignments,
all source completions and the relevant auxiliary extensions.  The audit also
validated the declared width-two decomposition, an independent-product
control, Horn and 2SAT controls, GF(2) Gaussian consistency controls, and the
wrapper family above.  Larger $n\in\{16,32,64,128,256,512\}$ used only the
proved closed form and deterministic parity samples; no asymptotic conclusion
is inferred from those samples.

\begin{center}
\begin{tabular}{lr}
\toprule
meter & value\\
\midrule
assertions & 40,582\\
source partial assignments & 39,360\\
compatibility evaluations & 88,147\\
completion checks & 118,180\\
CNF clause checks & 3,419,312\\
frontier entries explicitly examined & 1,020\\
preprocessing operations & 650,274\\
one-bit SAT decisions & 32\\
extraction / reconstruction / communication & 0 / 0 / 0\\
failures & 0\\
\bottomrule
\end{tabular}
\end{center}
At $n=8$ the exhaustive frontier size was 128 for each target and layout; at
$n=512$ the proved frontier size is $2^{511}$ while evaluation remains one
512-bit scan.  Script SHA-256 is
\texttt{f1925d9e1a4afba80029c43ffc21832e0e55dd4c666de2c8614bc74f725b7f09};
result SHA-256 is
\texttt{df830bb50ac1180bb4ae2d5abb83613f1409bd3f6d9c7fb60ed3336560843094}.

\section{Literature and barrier-transfer preflight}

The two blocking obstacles were: (i) turning a structural compatibility object
into an unavoidable event of every arbitrary computation, and (ii) landing
the resulting charge in one unconditional lower bound for the same ordinary
machine realization.  Primary literature was checked before the construction.

\begin{longtable}{p{0.20\linewidth}p{0.33\linewidth}p{0.38\linewidth}}
\toprule
source/domain & exact usable content & transfer classification\\
\midrule
Schaefer's Boolean CSP dichotomy & affine relations are polynomial-time
tractable; the present family is an explicit affine instance &
\textbf{directly transferable theorem}; it predicts the semantic shortcut,
not a hardness source\\
Tseitin transformations and parity formulas & linear-size gate encoding and
parity semantics & \textbf{directly transferable} after the four-clause gate
is proved here; proof-system hardness does not transfer to arbitrary SAT time\\
Samer--Szeider incidence-treewidth algorithms & SAT/\#SAT become tractable
under the declared decomposition parameter & \textbf{directly transferable}
to the structural parameter, but not needed for the sharper scan evaluator\\
Tseitin proof-complexity lower bounds & hardness for named proof systems &
\textbf{misleading analogy / failed transfer} to unrestricted computation\\
coding-theory Schur powers, relation spaces and syzygies & parity sources have
linear-code structure & \textbf{exact object correspondence}, but the
structure is algorithmically exploitable and supplies no runtime transfer\\
communication cuts, VLSI routing, observability and special soundness & lower
bounds in declared external/interface models & \textbf{useful analogy needing
a new transfer theorem}; arbitrary one-bit SAT exposes no such interface\\
\bottomrule
\end{longtable}

\paragraph{Post-quantum ledger.}
The source solution set is an affine code and its quadratic/Schur invariants
are exact algebraic objects, but Gaussian elimination and a parity accumulator
exploit them; classification: \textbf{FAILED TRANSFER} as hardness.  Collision
entropy, LWE-style worst/average-case architecture, and special-soundness
forking remain \textbf{ANALOGY/CANDIDATE} until a charged black-box extractor
from an arbitrary SAT decider is proved.  Hash/random-oracle binding is not an
unconditional theorem and is a \textbf{FAILED TRANSFER} for P versus NP.

\section{Barrier audit}

\paragraph{Relativization.}
The family theorem and recognize-first wrapper continue to hold relative to an
arbitrary oracle.  They cannot be the nonrelativizing hinge needed for a final
P-versus-NP proof.

\paragraph{Algebrization.}
The construction is affine and becomes easier with algebraic access.  It gives
no nonalgebrizing step.

\paragraph{Natural Proofs.}
No large constructive circuit property or unrestricted circuit lower bound is
proved, so the Natural Proofs barrier is not triggered or escaped.

\paragraph{Nonuniformity and arbitrary exponents.}
The evaluator is uniform.  Nothing bounds the unknown exponent of a hypothetical
full SAT decider, and no compiler maps it into a cited forbidden time--space
region.  Any claim suppressing that exponent would be circular.

\paragraph{\same.}
Running the displayed compatibility evaluator is a new computation chosen by
the analyst.  It is not a trace fact about an arbitrary execution of $A(F)$.
Conversely, counting semantic frontier members does not charge the one-bit run.

\section{Corpus reviewed}

\textbf{CORPUS REVIEWED.}  The full repository tree was surveyed.  Root
\path{main.tex}, \path{main4.tex}, \path{main5.tex}, \path{AGENTS.md},
\path{CODEX.md}, the complete pnp-research skill and all current
\path{tasks/} were read.  Materially relevant artifacts included newest MAIN
SYNTHESIS-0022; HEVALLB1; predecessor TESTS-0021; APIFLAT1; COREEQUIV1;
QGUARD1; VIEWSEP1; LB-0022; TESTS-0018/0020; the HEC, VSC-HEC, G-PPC-HEC,
CRS-PPC, CCSPC and RCCN causal branch; governance notes; and cited experiment
artifacts.  Repository search covered earlier ``minimal incompatible'',
``frontier'', ``partial assignment'', parity and compatibility formulations.

\textbf{PREDECESSOR WRITE-RECOVERY PREFLIGHT.}  The newest verified LSPV-2
artifact remains TESTS-0021 at 04:27 UTC, with valid header and sentinel.  No
substantive LSPV-2 artifact for the next scheduled cycle was present when this
run began.  The failure is classified \textbf{UNKNOWN}: there is no malformed
or partial target path and no connector error body from which to diagnose a
SHA conflict.  A body-preserving retry request is included in the handoff; no
scientific content was reconstructed.

\textbf{DUPLICATION / CONTRADICTION AUDIT.}  LB-0022 already showed exponentially
many full compatibility rows with one streaming parity bit.  HEVALLB1 showed
that a general full compatibility evaluator is polynomial exactly if SAT is.
VIEWSEP1 permitted compact separating evaluators.  The new contribution is
narrower: an exact $2^{n-1}$ antichain of \emph{minimal incompatible projected
partial assignments}, two concrete 3-CNF layouts, and a direct scan evaluator.
It strengthens the negative control but does not supersede QGUARD1.

\textbf{WHAT CHANGED RELATIVE TO ROOT MAIN.}  One stronger semantic-count
promotion is retired: minimality and antichain structure do not repair the
state-count-to-runtime gap.  The full-domain transfer and time--space landing
remain unchanged and OPEN.

\textbf{ADOPTED / REJECTED / RETIRED PRIOR COMPONENTS.}  Adopted: QGUARD1's
quantifier discipline, TESTS-0021's resource ledger, APIFLAT1's wrapper/API
guard, VIEWSEP1's compact-evaluator allowance, and COREEQUIV1's exact time-axis
boundary.  Rejected as transfers: proof-system Tseitin hardness, raw Schur or
Hilbert invariant size, and interface-model lower bounds without a compiler.
Retired: MIFO1 and any claim that making incompatible states minimal forces
their enumeration.

\textbf{STALE OR OVERSTATED CLAIMS FOUND.}  Any project prose that promotes an
exponential compatibility frontier directly to arbitrary computation cost is
overstated unless it names a full-domain execution relation and proves bounded
aggregation for the same realization.  No verified theorem establishing that
promotion was found.

\section{Status and handoff}

\begin{center}
\begin{tabular}{p{0.50\linewidth}p{0.28\linewidth}}
\toprule
claim & status\\
\midrule
PARITY-PROJECTED-FRONTIER1 & \proved, pending review\\
finite implementation and hostile-control audit & \comp\\
MINIMAL-INCOMPATIBILITY-FRONTIER-OBLIGATION1 & \retired\\
full-domain worst-case compatibility/execution transfer & \OPEN\\
model-matched nonrelativizing contradiction & \OPEN\\
P versus NP & \OPEN\\
\bottomrule
\end{tabular}
\end{center}

\paragraph{Smallest surviving repair.}
Prove a representation-invariant external relation $R(A,F,y)$ with the exact
prefix
\[
\forall A\in\mathcal D_{\mathrm{poly}}\ \exists F\in X_N\ \forall N\ge N_0(A):
\quad R(A,F,y)\text{ forces a charged event of the same execution},
\]
or correct that prefix explicitly, then show bounded aggregation and reduce the
resulting charge to one cited unconditional lower bound for the same uniform
random-access machine.  Easy islands, proof objects and analyst-run evaluators
cannot supply this arrow.

\paragraph{What this does not prove.}
It proves neither a general SAT algorithm nor a SAT lower bound; it does not
show that every hard family has compact compatibility; it does not refute a
properly quantified worst-case transfer; and it does not resolve P versus NP.

\begin{thebibliography}{10}
\bibitem{Schaefer1978}
T.~J. Schaefer, ``The complexity of satisfiability problems,''
\emph{Proceedings of STOC}, 1978, pp.~216--226.

\bibitem{Tseitin1968}
G.~S. Tseitin, ``On the complexity of derivation in propositional calculus,''
in \emph{Studies in Constructive Mathematics and Mathematical Logic}, 1968.

\bibitem{SamerSzeider2010}
M.~Samer and S.~Szeider, ``Algorithms for propositional model counting,''
\emph{Journal of Discrete Algorithms} 8(1), 2010, pp.~50--64,
\url{https://arxiv.org/abs/0610174}.

\bibitem{BGS1975}
T.~Baker, J.~Gill and R.~Solovay, ``Relativizations of the $P=?NP$ question,''
\emph{SIAM Journal on Computing} 4(4), 1975, pp.~431--442.

\bibitem{AW2009}
S.~Aaronson and A.~Wigderson, ``Algebrization: a new barrier in complexity
theory,'' \emph{ACM Transactions on Computation Theory} 1(1), 2009.

\bibitem{RR1997}
A.~Razborov and S.~Rudich, ``Natural proofs,'' \emph{Journal of Computer and
System Sciences} 55(1), 1997, pp.~24--35.
\end{thebibliography}

% END_OF_VERIFIED_REPORT::LSPV-3::MINFRONT1
\end{document}
